\setcounter{page}{100}
\section{\S5 Relations Among the Derived Functors}

In this section we list natural homomorphisms and isomorphisms between the derived functors defined in the previous sections.

\textbf{Proposition 5.1.}
Let $X\xrightarrow{f} Y\xrightarrow{g} Z$ be morphisms of preschemes.
There is a natural isomorphism
\[
\zeta^+\colon \mathbf{R}^+(g_*\circ f_*) \xrightarrow{\sim} \mathbf{R}^+g_*\circ\mathbf{R}^+f_*
\]
of functors $D^+(X)\to D^+(Z)$.

If furthermore $X,Y$ are noetherian of finite Krull dimension, there is a natural isomorphism
\[
\zeta\colon \mathbf{R}(g_*\circ f_*) \xrightarrow{\sim} \mathbf{R}g_*\circ\mathbf{R}f_*
\]
of functors $D(X)\to D(Z)$.

\textit{Proof.}
Use Theorem I.5.4(b).
For the bounded case, let $L\subseteq K^+(\operatorname{Mod}(X))$ be complexes of injective sheaves,
and $M\subseteq K^+(\operatorname{Mod}(Y))$ complexes of $g_*$-acyclic sheaves.
Then $f_*(L)\subseteq M$: injective sheaves are flasque (EGA $\mathrm{O_{III}}$ 12.2.4),
and flasque sheaves are $g_*$-acyclic (EGA $\mathrm{O_{III}}$ 12.2.1).
The universal property yields $\zeta^+$.

\setcounter{page}{101}
For the unbounded case, let $L\subseteq K(\operatorname{Mod}(X))$ be complexes that are both $f_*$-acyclic and $g_*f_*$-acyclic,
and $M\subseteq K(\operatorname{Mod}(Y))$ complexes of $g_*$-acyclic sheaves.
If $F\in L$, then $\mathbf{R}f_*(F)\cong f_*(F)$,
so $R^ig_*(f_*F)=R^i(g_*f_*)(F)=0$ for $i>0$, hence $f_*F\in M$.
The class of sheaves which are $f_*$- and $g_*f_*$-acyclic satisfies I.4.6,
so Theorem I.5.4(b) applies to give the isomorphism $\zeta$.

\textbf{Proposition 5.2.}
Let $f\colon X\to Y$ be a morphism of preschemes.
There is a natural isomorphism
\[
\zeta^+\colon \mathbf{R}^+\Gamma(X,\,\cdot\,) \xrightarrow{\sim} \mathbf{R}^+\Gamma(Y,\mathbf{R}^+f_*)
\]
of functors $D^+(X)\to D(\mathbf{Ab})$.

If $X,Y$ are noetherian of finite Krull dimension, there is a natural isomorphism
\[
\zeta\colon \mathbf{R}\Gamma(X,\,\cdot\,) \xrightarrow{\sim} \mathbf{R}\Gamma(Y,\mathbf{R}f_*)
\]
of functors $D(X)\to D(\mathbf{Ab})$.

\setcounter{page}{102}
\textit{Proof.}
Similar to Proposition 5.1.

\textbf{Proposition 5.3.}
Let $X$ be a prescheme. There is a natural isomorphism of bifunctors
\[
\mathbf{R}\mathcal{H}\textit{om}^\bullet(F^\bullet,G^\bullet)
\xrightarrow{\sim}
\Gamma\big(X,\mathbf{R}\underline{\mathcal{H}\textit{om}}^\bullet(F^\bullet,G^\bullet)\big)
\]
from $D^-(X)^\circ\times D^+(X)$ to $D(\mathbf{Ab})$.

\textit{Proof.}
For ordinary sheaves, $\Hom(F,G)=\Gamma(X,\underline{\Hom}(F,G))$.
If $G$ is injective, $\underline{\Hom}(F,G)$ is flasque [G, II 7.3.2], hence $\Gamma$-acyclic.
The result follows from derived functor universal properties.

\textbf{Proposition 5.4.}
Let $X\xrightarrow{f} Y\xrightarrow{g} Z$ be morphisms of preschemes.
There is a natural isomorphism
\[
\zeta^-\colon \mathbf{L}^-(f^*\circ g^*)
\xrightarrow{\sim}
\mathbf{L}^-f^*\circ\mathbf{L}^-g^*
\]
of functors $D^-(Z)\to D^-(X)$.

If $f,g$ have finite Tor-dimension, then so does $g\circ f$, and there is a natural isomorphism
\[
\zeta\colon \mathbf{L}(g\circ f)^*
\xrightarrow{\sim}
\mathbf{L}f^*\circ\mathbf{L}g^*
\]
of functors $D(Z)\to D(X)$.

\setcounter{page}{103}
\textit{Proof.}
Left to the reader. Note $g^*$ sends flat $\sheaf{O}_Z$-modules to flat $\sheaf{O}_Y$-modules.

\textbf{Proposition 5.5.}
Let $f\colon X\to Y$, with $X$ noetherian of finite Krull dimension.
For $F^\bullet\in D^-(X)$, $G^\bullet\in D^+(X)$, there is a natural functorial homomorphism
\[
\mathbf{R}f_*\mathbf{R}\mathcal{H}\textit{om}_X(F^\bullet,G^\bullet)
\to
\mathbf{R}\mathcal{H}\textit{om}_Y(\mathbf{R}f_*F^\bullet,\mathbf{R}f_*G^\bullet).
\]

\textit{Proof.}
The hypotheses guarantee both sides are well-defined (\S3).
We define a morphism of bifunctors $D^-(X)^\circ\times D^+(X)\to D(Y)$.
Let $L\subseteq K^-(\operatorname{Mod}(X))$ be complexes of $f_*$-acyclic sheaves,
$M\subseteq K^+(\operatorname{Mod}(X))$ complexes of injective sheaves.
The localizations $L_{\mathsf{Qis}}\to D^-(X)$ and $M_{\mathsf{Qis}}\to D^+(X)$ are equivalences,
so it suffices to define the morphism on $L_{\mathsf{Qis}}\times M_{\mathsf{Qis}}$.

\setcounter{page}{104}
Using the canonical morphisms from derived functor definitions, we have the commutative diagram:
\[
\begin{array}{l}
(1)\ \xi_{f_*}\colon \mathbf{R}f_*\underline{\mathcal{H}\textit{om}}^\bullet(F^\bullet,G^\bullet)
\to \mathbf{R}f_*\mathbf{R}\underline{\mathcal{H}\textit{om}}^\bullet(F^\bullet,G^\bullet) \\
(2)\ \xi_{\mathcal{H}\textit{om}}\colon \underline{\mathcal{H}\textit{om}}^\bullet(F^\bullet,G^\bullet)
\to \mathbf{R}\underline{\mathcal{H}\textit{om}}^\bullet(F^\bullet,G^\bullet) \\
(4)\ \xi_{f_*}\colon \underline{\mathcal{H}\textit{om}}^\bullet(f_*F^\bullet,f_*G^\bullet)
\to \mathbf{R}\underline{\mathcal{H}\textit{om}}^\bullet(f_*F^\bullet,f_*G^\bullet) \\
(5)\ \xi_{f_*}\colon \mathbf{R}\mathcal{H}\textit{om}_Y(f_*F^\bullet,f_*G^\bullet)
\to \mathbf{R}\mathcal{H}\textit{om}_Y(\mathbf{R}f_*F^\bullet,\mathbf{R}f_*G^\bullet)
\end{array}
\]

The map is induced by the canonical sheaf morphism
\[
f_*\underline{\Hom}_X(F,G)
\to
\underline{\Hom}_Y(f_*F,f_*G).
\]

\setcounter{page}{105}
(1) is an isomorphism since $\underline{\mathcal{H}\textit{om}}^\bullet(F^\bullet,G^\bullet)$ is flasque hence $f_*$-acyclic.
(2) is an isomorphism as $G^\bullet$ is injective.
(4), (5) are likewise isomorphisms.
Thus there exists a unique functor morphism
\[
\Psi\colon
\mathbf{R}f_*\mathbf{R}\underline{\mathcal{H}\textit{om}}^\bullet(F^\bullet,G^\bullet)
\to
\mathbf{R}\mathcal{H}\textit{om}_Y(\mathbf{R}f_*F^\bullet,\mathbf{R}f_*G^\bullet)
\]
making the diagram commute.

\textbf{Remark.}
In subsequent arguments we adopt the convention:
for complexes of acyclic/injective objects, we identify $\mathbf{R}F$ with $F$ and omit localization functors $Q$.

\setcounter{page}{106}
\textbf{Proposition 5.6 (Projection Formula).}
Let $f\colon X\to Y$ be quasi-compact, with $X,Y$ noetherian of finite Krull dimension.
For $F^\bullet\in D^-(X)$, $G^\bullet\in D_{\mathrm{qc}}^-(Y)$, there is a natural isomorphism
\[
\mathbf{R}f_*\big(F^\bullet \stackrel{\mathbf{L}}{\otimes}_X \mathbf{L}f^*G^\bullet\big)
\;\simeq\;
\mathbf{R}f_*F^\bullet \stackrel{\mathbf{L}}{\otimes}_Y G^\bullet.
\]

\textit{Proof.}
Assume $F^\bullet$ is $f_*$-acyclic and $G^\bullet$ is a complex of flat $\sheaf{O}_Y$-modules.
Compose the classical sheaf projection formula
\[
f_*\big(F^\bullet\otimes_X f^*G^\bullet\big)
\to
f_*F^\bullet\otimes_Y G^\bullet
\]
with derived functor identifications to obtain the isomorphism.